% =============================================================================
%  CHAPTER 9: Code Listings & Verbatim
% =============================================================================

\chapter{Code Listings \& Verbatim Text}

\section{Verbatim Text}

The \texttt{verbatim} environment prints text exactly as typed, without interpreting any \latex{} commands.

\begin{lstlisting}[language={[LaTeX]TeX}, caption=Verbatim environments]
% Inline verbatim:
\verb|This is \verbatim{test} text|

% Display verbatim:
\begin{verbatim}
    \textbf{This text is NOT bold}
    $E = mc^2$ is NOT an equation
    Everything appears exactly as typed!
\end{verbatim}

% Verbatim with different delimiter:
\verb!Special chars: $ % & # _ { }!

% With fancyvrb package (more options):
\usepackage{fancyvrb}
\begin{Verbatim}[fontsize=\small, frame=single, framesep=2mm]
    Custom verbatim with frame and smaller font
\end{Verbatim}
\end{lstlisting}

\subsection{Inline Verbatim Examples}

\begin{itemize}
    \item \verb|\documentclass{article}| — the document class command
    \item \verb|\usepackage{graphicx}| — loading a package
    \item \verb|$E = mc^2$| — inline math
\end{itemize}

% ------------------------------------------------------------------
\section{Code Listings (listings package)}

The \texttt{listings} package provides sophisticated source code formatting.

\subsection{Basic Listing}

\begin{lstlisting}[language=Python, caption=Basic code listing]
def hello_world():
    print("Hello, World!")
\end{lstlisting}

\subsection{Listing with Language}

\begin{lstlisting}[language=Python, caption=Python example with syntax highlighting]
# Python example: Fibonacci sequence
def fibonacci(n):
    """Calculate the nth Fibonacci number."""
    if n <= 1:
        return n
    a, b = 0, 1
    for _ in range(2, n + 1):
        a, b = b, a + b
    return b

# Generate first 10 Fibonacci numbers
for i in range(10):
    print(f"F({i}) = {fibonacci(i)}")
\end{lstlisting}

\subsection{Java Example}

\begin{lstlisting}[language=Java, caption=Java example]
public class HelloWorld {
    public static void main(String[] args) {
        System.out.println("Hello, World!");

        // Loop example
        for (int i = 0; i < 5; i++) {
            System.out.println("Iteration: " + i);
        }
    }
}
\end{lstlisting}

\subsection{C/C++ Example}

\begin{lstlisting}[language=C++, caption=C++ example]
#include <iostream>
#include <vector>
#include <algorithm>

int main() {
    std::vector<int> numbers = {5, 2, 8, 1, 9, 3};

    // Sort ascending
    std::sort(numbers.begin(), numbers.end());

    // Print sorted vector
    for (const auto& n : numbers) {
        std::cout << n << " ";
    }
    std::cout << std::endl;

    return 0;
}
\end{lstlisting}

\subsection{JavaScript Example}

\begin{lstlisting}[language=Java, caption=JavaScript example]
// Asynchronous data fetching
async function fetchData(url) {
    try {
        const response = await fetch(url);
        const data = await response.json();
        return data;
    } catch (error) {
        console.error("Error:", error);
    }
}
\end{lstlisting}

% ------------------------------------------------------------------
\section{Listing Styles and Configuration}

\begin{lstlisting}[language={[LaTeX]TeX}, caption=Listing configuration]
% Global settings in preamble:
\lstset{
    basicstyle=\ttfamily\small,          % Base font
    keywordstyle=\color{blue}\bfseries,  % Keywords
    commentstyle=\color{green!50!black}\itshape, % Comments
    stringstyle=\color{red!70!black},    % Strings
    numbers=left,                        % Line numbers
    numberstyle=\tiny\color{gray},       % Line number style
    stepnumber=1,                        % Line number step
    numbersep=5pt,                       % Distance of numbers
    backgroundcolor=\color{gray!5},      % Background color
    frame=single,                        % Border frame
    rulecolor=\color{gray!50},           % Frame color
    breaklines=true,                     % Auto line break
    breakatwhitespace=false,             % Break only at whitespace
    tabsize=4,                           % Tab size
    showstringspaces=false,              % Don't show string spaces
    captionpos=b,                        % Caption at bottom
    title=\lstname,                      % Show filename as title
}
\end{lstlisting}

\note{Per-listing options go in square brackets: \texttt{\textbackslash begin\{lstlisting\}[language=Python, caption=\{...\}]}.}

% ------------------------------------------------------------------
\section{Inline Code}

\begin{lstlisting}[language={[LaTeX]TeX}, caption=Inline code]
% Use \lstinline for inline code:
The function \lstinline|print("hello")| outputs text.

% Short form (define a shorthand):
\lstMakeShortInline[language=Python]{|}
Use |print("hello")| in your code.
\end{lstlisting}

% ------------------------------------------------------------------
\section{Supported Languages}

The listings package supports many programming languages:

\begin{table}[H]
\centering
\begin{tabular}{llll}
\toprule
\textbf{Language} & \textbf{Keyword} & \textbf{Language} & \textbf{Keyword} \\
\midrule
Python     & \texttt{Python}        & Java      & \texttt{Java} \\
C          & \texttt{C}             & C++       & \texttt{C++} \\
JavaScript & \texttt{Java}          & HTML      & \texttt{HTML} \\
SQL        & \texttt{SQL}           & R         & \texttt{R} \\
MATLAB     & \texttt{Matlab}        & Fortran   & \texttt{Fortran} \\
Bash       & \texttt{bash}          & XML       & \texttt{XML} \\
LaTeX      & \texttt{[LaTeX]TeX}    & Make      & \texttt{make} \\
\bottomrule
\end{tabular}
\caption{Some supported languages in the listings package}
\end{table}
